\documentclass[12pt,a4paper]{article}

\usepackage[utf8]{inputenc}
\usepackage[T1]{fontenc}
\usepackage[spanish]{babel}
\usepackage{amsmath, amssymb}
\usepackage{geometry}
\usepackage{booktabs}
\usepackage{array}
\usepackage{xcolor}
\usepackage{hyperref}
\usepackage{enumitem}
\usepackage{fancyhdr}
\usepackage{graphicx}
\usepackage{listings}
\usepackage{tcolorbox}
\usepackage{multirow}
\usepackage{caption}

\geometry{top=2.5cm, bottom=2.5cm, left=3cm, right=3cm}

% ── Colores ──────────────────────────────────────────────────────────────────
\definecolor{uddblue}{RGB}{0, 75, 135}
\definecolor{lightgray}{RGB}{245, 245, 245}
\definecolor{codegray}{RGB}{100, 100, 100}
\definecolor{codegreen}{RGB}{0, 120, 0}
\definecolor{codeblue}{RGB}{0, 0, 180}
\definecolor{codered}{RGB}{160, 0, 0}

% ── Estilo de código Python ───────────────────────────────────────────────────
\lstset{
    language     = Python,
    backgroundcolor = \color{lightgray},
    basicstyle   = \small\ttfamily,
    keywordstyle = \color{codeblue}\bfseries,
    commentstyle = \color{codegreen}\itshape,
    stringstyle  = \color{codered},
    numberstyle  = \tiny\color{codegray},
    numbers      = left,
    stepnumber   = 1,
    breaklines   = true,
    frame        = single,
    rulecolor    = \color{gray!40},
    tabsize      = 4,
    showstringspaces = false,
    literate =
        {á}{{\'{a}}}1 {é}{{\'{e}}}1 {í}{{\'{i}}}1
        {ó}{{\'{o}}}1 {ú}{{\'{u}}}1 {ñ}{{\~{n}}}1
        {Á}{{\'{A}}}1 {É}{{\'{E}}}1 {Í}{{\'{I}}}1
        {Ó}{{\'{O}}}1 {Ú}{{\'{U}}}1 {Ñ}{{\~{N}}}1
}

% ── Cajas de ayuda ────────────────────────────────────────────────────────────
\tcbuselibrary{skins, breakable}
\newtcolorbox{ayuda}[1]{
    colback   = lightgray,
    colframe  = uddblue,
    fonttitle = \bfseries\small,
    title     = {#1},
    breakable,
    left      = 6pt,
    right     = 6pt,
    top       = 4pt,
    bottom    = 4pt
}

% ── Encabezado y pie ─────────────────────────────────────────────────────────
\pagestyle{fancy}
\fancyhf{}
\fancyhead[L]{\small\color{gray} IIP314W\,--\,Optimización Aplicada a Negocios}
\fancyhead[R]{\small\color{gray} Tarea 1\,--\,2026-T2}
\fancyfoot[C]{\small\thepage}
\renewcommand{\headrulewidth}{0.4pt}

\hypersetup{colorlinks=true, linkcolor=uddblue, urlcolor=uddblue}

% ── Espaciado entre párrafos ──────────────────────────────────────────────────
\setlength{\parskip}{0.5em}
\setlength{\parindent}{0pt}

% =============================================================================
\begin{document}
% =============================================================================

% ── Título ───────────────────────────────────────────────────────────────────
\begin{center}
    {\Large\bfseries\color{uddblue}
     (IIP314W) Optimización Aplicada a Negocios}\\[0.4em]
    {\large Universidad del Desarrollo}\\[0.8em]
    \noindent\rule{\linewidth}{1.5pt}\\[0.5em]
    {\LARGE\bfseries Tarea 1}\\[0.3em]
    {\large Métodos de Optimización Numérica:\\
     Gradiente Descendente y Newton-Raphson}\\[0.5em]
    \noindent\rule{\linewidth}{1.5pt}\\[0.8em]
    \begin{tabular}{@{}ll@{}}
        \textbf{Profesor:} & Ing.\ Rodrigo Trigo Vilches \\[0.3em]
        \textbf{Ayudante:} & Lic.\ Vicente Ramírez Almonacid \\[0.3em]
        \textbf{Período:}  & Año 2026\,--\,Trimestre 2
    \end{tabular}
\end{center}

\vspace{1em}
\noindent\rule{\linewidth}{0.4pt}
\vspace{0.5em}

% =============================================================================
\section*{Contexto}
% =============================================================================

En muchos problemas de negocio —minimización de costos, ajuste de modelos
predictivos, calibración de parámetros— la función que se desea optimizar
\textbf{no es convexa}: presenta múltiples mínimos locales. En estos casos, el
algoritmo de optimización y, sobre todo, el \textbf{punto de partida} elegido
determinan a qué solución se converge. Comprender este comportamiento es
fundamental antes de aplicar métodos numéricos a problemas reales.

En esta tarea estudiarán y compararán tres métodos clásicos de optimización
numérica sobre una misma función de prueba multimodal:

\begin{enumerate}[label=\arabic*., leftmargin=*, topsep=2pt]
    \item Gradiente descendente con \textbf{paso fijo}.
    \item Gradiente descendente con \textbf{dirección normalizada}.
    \item Método de \textbf{Newton-Raphson}.
\end{enumerate}

Considere la siguiente función de dos variables, compuesta por la suma de
tres gaussianas: dos \emph{invertidas} (negativas) y una \emph{positiva}, cada
una centrada en un punto distinto del plano. Esto genera \textbf{dos mínimos
locales y un máximo local} bien diferenciados:

\begin{equation*}
    f(x, y) = -\,e^{-\left[(x+2)^2+(y+1)^2\right]}
              \;-\; 1{,}3\,e^{-\left[(x-2)^2+(y+1)^2\right]}
              \;+\; 1{,}5\,e^{-\left[x^2+(y-1{,}5)^2\right]}
\end{equation*}

Esta estructura es un excelente caso de estudio: un método que solo sabe
``bajar'' (gradiente descendente) se comportará de forma muy distinta a uno
que busca anular el gradiente (Newton-Raphson), el cual no distingue entre
mínimos, máximos y puntos silla. Le anticipamos, además, que entre ambos
mínimos aparece de forma natural un \textbf{punto silla}: identificarlo y
clasificarlo es parte del análisis.

\begin{ayuda}{Trabajo en grupo}
\vspace{2pt}
La tarea se desarrolla \textbf{en parejas o grupos de 3} (ni más ni menos).
Tanto el informe como el código deben incluir los nombres de todos los
integrantes.
\end{ayuda}

\noindent\rule{\linewidth}{0.4pt}

% =============================================================================
\section*{Parte A \quad Análisis Gráfico \hfill\normalfont\normalsize(10 pts)}
% =============================================================================

\begin{enumerate}[label=\alph*), leftmargin=*, topsep=2pt]
    \item Visualice la función en una \textbf{gráfica 3D} para identificar la
          estructura de sus mínimos y de su máximo.
    \item Genere un \textbf{mapa de curvas de nivel} (contour) que permita
          apreciar los valles, el cerro (máximo) y las regiones planas de la
          función.
    \item A partir de los gráficos, estime visualmente la ubicación aproximada
          de los \textbf{dos mínimos} y del \textbf{máximo}. Estas estimaciones
          le servirán para elegir puntos iniciales y para validar los
          resultados numéricos.
\end{enumerate}

\begin{ayuda}{Pista -- Herramientas de visualización}
\vspace{2pt}
Para la superficie 3D necesitará una malla de puntos (\texttt{numpy.meshgrid})
y un eje con proyección \texttt{"3d"} (método \texttt{plot\_surface}). Para las
curvas de nivel, las funciones \texttt{contour} / \texttt{contourf} de
\texttt{matplotlib} reciben esa misma malla. Si algún punto crítico queda fuera
de la ventana, ajuste los rangos de $x$ e $y$ (sugerencia: explore en torno a
$x\in[-5,5]$, $y\in[-4,5]$).
\end{ayuda}

\noindent\rule{\linewidth}{0.4pt}

% =============================================================================
\section*{Parte B \quad Gradiente Descendente con Paso Fijo \hfill\normalfont\normalsize(20 pts)}
% =============================================================================

\begin{enumerate}[label=\alph*), leftmargin=*, topsep=2pt]
    \item Implemente el algoritmo de gradiente descendente con un
          \textbf{tamaño de paso fijo} $\alpha$.
    \item Pruebe con \textbf{al menos 4 puntos iniciales} distintos,
          seleccionados estratégicamente en distintas regiones del dominio
          (cercanos a cada valle, cerca del máximo y en zonas intermedias).
    \item Para cada punto inicial, registre el \textbf{mínimo local encontrado},
          el \textbf{valor de la función} en ese punto y el \textbf{número de
          iteraciones} hasta la convergencia.
    \item Visualice las \textbf{trayectorias de convergencia} sobre el mapa de
          curvas de nivel.
\end{enumerate}

\begin{ayuda}{Pista -- Regla de actualización}
\vspace{2pt}
La iteración del gradiente descendente con paso fijo es:
\[
    \mathbf{x}_{k+1} = \mathbf{x}_k - \alpha \,\nabla f(\mathbf{x}_k),
\]
donde $\alpha > 0$ es el paso (tasa de aprendizaje) y $\nabla f$ es el
gradiente. Un criterio de parada habitual es
$\lVert \nabla f(\mathbf{x}_k) \rVert < \texttt{tol}$ o alcanzar un número
máximo de iteraciones. Experimente con el valor de $\alpha$: un paso muy
grande puede hacer divergir el método, y uno muy pequeño puede volverlo
demasiado lento.
\end{ayuda}

\noindent\rule{\linewidth}{0.4pt}

% =============================================================================
\section*{Parte C \quad Gradiente Descendente con Dirección Normalizada \hfill\normalfont\normalsize(20 pts)}
% =============================================================================

\begin{enumerate}[label=\alph*), leftmargin=*, topsep=2pt]
    \item Implemente el algoritmo de gradiente descendente con
          \textbf{dirección normalizada}: en cada paso se avanza una distancia
          fija en la dirección del gradiente, independientemente de su
          magnitud.
    \item Use los \textbf{mismos puntos iniciales} que en la Parte~B.
    \item Registre el mínimo local encontrado, el valor de la función y el
          número de iteraciones en cada caso.
    \item Visualice las trayectorias de convergencia sobre el mapa de curvas
          de nivel.
\end{enumerate}

\begin{ayuda}{Pista -- Dirección normalizada}
\vspace{2pt}
La regla de actualización usa el gradiente \emph{normalizado}:
\[
    \mathbf{x}_{k+1} = \mathbf{x}_k
    - \alpha \,\frac{\nabla f(\mathbf{x}_k)}{\lVert \nabla f(\mathbf{x}_k)\rVert}.
\]
Así, el tamaño del paso es siempre $\alpha$ (en distancia), sin importar si el
gradiente es grande o pequeño. Reflexione sobre cómo afecta esto a la
convergencia cerca del mínimo, donde el gradiente tiende a cero.
\end{ayuda}

\noindent\rule{\linewidth}{0.4pt}

% =============================================================================
\section*{Parte D \quad Método de Newton-Raphson \hfill\normalfont\normalsize(25 pts)}
% =============================================================================

\begin{enumerate}[label=\alph*), leftmargin=*, topsep=2pt]
    \item Investigue el método de Newton-Raphson para encontrar
          \textbf{raíces de un sistema de ecuaciones} y aplíquelo para hallar
          los \textbf{puntos críticos} de $f$ (es decir, los puntos donde
          $\nabla f = \mathbf{0}$).
    \item Para ello necesitará el \textbf{gradiente} $\nabla f$ y la
          \textbf{matriz Hessiana} $H_f$. Puede obtenerlos de forma analítica
          (a mano o con software simbólico) o numérica; en cualquier caso,
          debe quedar documentado en el informe.
    \item Al igual que en los métodos anteriores, use los \textbf{mismos
          puntos iniciales}. Registre el punto crítico encontrado, el valor de
          la función y el número de iteraciones.
    \item Visualice las trayectorias de convergencia sobre el mapa de curvas
          de nivel. \textbf{Comente} si el método siempre converge a un mínimo
          o si, dependiendo del punto inicial, puede converger a un máximo o a
          un punto silla.
\end{enumerate}

\begin{ayuda}{Pista -- Este método deben investigarlo}
\vspace{2pt}
A diferencia de los métodos anteriores, \textbf{aquí no se entrega la regla de
actualización}: forma parte del trabajo investigar el método de Newton-Raphson
y deducir cómo aplicarlo a este problema. Algunas preguntas que los pueden
guiar:
\begin{itemize}[topsep=2pt]
    \item ¿Qué condición matemática cumple \emph{cualquier} punto crítico
          (mínimo, máximo o silla)?
    \item Newton-Raphson resuelve sistemas de ecuaciones no lineales
          $\mathbf{F}(\mathbf{x}) = \mathbf{0}$. ¿Qué deberían tomar como
          $\mathbf{F}$ para hallar los puntos críticos de $f$?
    \item Una vez planteado así, ¿qué objeto matemático cumple el rol que la
          derivada $f'$ cumple en el Newton-Raphson escalar de una variable?
\end{itemize}
Citen la(s) fuente(s) que consulten en el informe.
\end{ayuda}

\noindent\rule{\linewidth}{0.4pt}
\newpage
% =============================================================================
\section*{Parte E \quad Análisis Comparativo \hfill\normalfont\normalsize(25 pts)}
% =============================================================================

\begin{enumerate}[label=\alph*), leftmargin=*, topsep=2pt]
    \item En un \textbf{gráfico comparativo} (lado a lado o en vertical),
          muestre las trayectorias de los tres métodos partiendo desde los
          mismos puntos iniciales sobre el mapa de curvas de nivel.
    \item Construya una \textbf{tabla resumen} con los resultados de cada
          método para cada punto inicial. Puede usar el siguiente formato
          (agregue columnas si lo estima conveniente):

          \begin{center}
          \renewcommand{\arraystretch}{1.2}
          \begin{tabular}{lccccc}
          \toprule
          \textbf{Método} &
          \textbf{Punto inicial} &
          \textbf{Punto encontrado} &
          \textbf{Punto crítico real} &
          \textbf{Gap} &
          \textbf{Iteraciones} \\
          \midrule
          GD paso fijo & $(x_0,y_0)$ & $(x_a,y_a)$ & $(x_b,y_b)$ & $m\%$ & $n$ \\
          \multicolumn{6}{c}{\dots} \\
          \bottomrule
          \end{tabular}
          \captionof{table}{Formato sugerido para la tabla comparativa.
          Fuente: elaboración propia.}
          \end{center}

    \item Analice las diferencias entre los tres métodos en términos de:
          \begin{itemize}[topsep=2pt]
              \item \textbf{Número de iteraciones} hasta la convergencia.
              \item \textbf{Sensibilidad al tamaño del paso} $\alpha$.
              \item \textbf{Robustez} ante distintos puntos iniciales.
              \item \textbf{Gap} entre los puntos críticos \emph{reales} y los
                    efectivamente encontrados.
              \item \textbf{Tipo de punto crítico} alcanzado por cada método:
                    ¿a qué punto(s) puede converger Newton que GD nunca
                    alcanza? Interprete la diferencia.
          \end{itemize}
\end{enumerate}

\begin{ayuda}{Pista -- Puntos críticos ``reales'' de referencia}
\vspace{2pt}
Para calcular el \textbf{Gap} necesita los puntos críticos reales (los valores
exactos hacia los que \emph{debería} converger cada método). Puede obtenerlos
resolviendo $\nabla f = \mathbf{0}$ a mano, con software simbólico, o con una
herramienta numérica de su elección; justifique en el informe cómo lo hizo.
El \textbf{Gap} puede expresarse, por ejemplo, como la distancia entre el punto
encontrado y el punto crítico real, o como la diferencia relativa en el valor
de la función. Defina y justifique la métrica que utilice.
\end{ayuda}

\noindent\rule{\linewidth}{0.4pt}

% =============================================================================
\section*{Entregables}
% =============================================================================

\subsection*{1.\quad Informe (PDF)}

El informe es la pieza central de la evaluación. Debe contener, como
\textbf{mínimo}, las siguientes secciones (cada una comenzando en una nueva
página):

\begin{enumerate}[leftmargin=*, topsep=2pt]
    \item \textbf{Portada:} nombre del curso, Tarea~1, integrantes del grupo,
          profesor, ayudante y fecha.
    \item \textbf{Introducción al problema:} descripción del problema de
          optimización abordado y de la función de estudio. ¿Por qué es
          relevante analizar el comportamiento de los métodos ante una función
          multimodal?
    \item \textbf{Marco teórico:} definición de los conceptos clave utilizados
          —gradiente, Hessiano, tasa de aprendizaje (paso), dirección de
          descenso, criterio de tolerancia, punto crítico, convergencia, etc.
    \item \textbf{Objetivos:} un \textbf{objetivo general} y al menos
          \textbf{tres objetivos específicos} que guíen el desarrollo.
    \item \textbf{Desarrollo (metodología):} explicación de \emph{cómo} se
          realizó la tarea de inicio a fin —cómo se obtuvieron el gradiente y
          el Hessiano, cómo se implementó cada método, qué puntos iniciales y
          qué parámetros ($\alpha$, tolerancia, máximo de iteraciones) se
          eligieron y por qué.
    \item \textbf{Resultados:} presente \textbf{solo los resultados
          importantes} —tabla comparativa, gráficos de trayectorias, gráfico
          comparativo de los tres métodos.
          \begin{quote}
          \textit{Ojo:} en esta sección van los resultados puros, sin
          interpretación.
          \end{quote}
    \item \textbf{Análisis de resultados:} interpretación de los resultados
          anteriores. Discuta las diferencias de iteraciones, la sensibilidad
          al paso, la robustez ante el punto inicial y los Gap observados.
    \item \textbf{Conclusiones:} cierre respecto a la importancia de los
          parámetros y del punto inicial, y si se cumplieron los objetivos
          planteados. Sea \textbf{específico}: evite conclusiones genéricas del
          tipo ``los parámetros son muy importantes'' o ``los métodos
          funcionan''.
    \item \textbf{Referencias} (si aplica).
\end{enumerate}

\begin{ayuda}{Formalidades del informe (se descuentan puntos por incumplimiento)}
\begin{itemize}[topsep=2pt]
    \item Lenguaje técnico y formal, sin errores ortográficos ni gramaticales.
    \item Texto \textbf{justificado}, fuente tamaño \textbf{12}
          (Times New Roman, Arial u otra tipografía estándar).
    \item Cada figura, gráfico o tabla debe llevar \textbf{título y fuente}.
          Si el material es de elaboración propia, debe indicarse
          explícitamente: \emph{Fuente: elaboración propia}. Si proviene de una
          fuente externa, debe citarse.
    \item Sean \textbf{breves y directos}. Se prefiere un análisis corto,
          directo y correcto, por sobre páginas de relleno.
\end{itemize}
\end{ayuda}

\subsection*{2.\quad Cuaderno Jupyter (.ipynb)}

\begin{itemize}[topsep=2pt]
    \item Implementación \textbf{comentada} de los tres métodos
          (Partes~B, C y D).
    \item Los gráficos solicitados (superficie 3D, curvas de nivel,
          trayectorias y gráfico comparativo).
    \item La tabla resumen con los resultados para cada punto inicial.
    \item Código limpio, ordenado y bien estructurado.
\end{itemize}

\medskip
\noindent\textbf{Formato de entrega:} comprima el informe y el notebook en un
archivo \texttt{.zip} o \texttt{.rar} con el nombre \texttt{Tarea1\_GrupoN}.
\textbf{No entregue los archivos por separado.}

\noindent\rule{\linewidth}{0.4pt}

% =============================================================================
\section*{Plazos y Evaluación}
% =============================================================================

\begin{itemize}[topsep=2pt]
    \item \textbf{Plazo de entrega:} \emph{por confirmar} (se anunciará en
          Canvas).
    \item \textbf{Grupos:} 2 a 3 integrantes (ni más ni menos).
    \item \textbf{Inscripción:} anótense en los grupos de Canvas
          (Personas~$\to$ Grupos~$\to$ Tarea~1\,--\,Grupo~n).
\end{itemize}

\vspace{0.5em}

\begin{center}
\renewcommand{\arraystretch}{1.2}
\begin{tabular}{lc}
\toprule
\textbf{Criterio} & \textbf{Ponderación} \\
\midrule
Parte A: Análisis gráfico                              & 10\,\% \\
Parte B: Gradiente descendente con paso fijo           & 20\,\% \\
Parte C: Gradiente descendente con dirección normalizada & 20\,\% \\
Parte D: Método de Newton-Raphson                      & 25\,\% \\
Parte E: Análisis comparativo                          & 25\,\% \\
\midrule
\textit{Calidad del informe y profundidad del análisis} &
\textit{transversal} \\
\bottomrule
\end{tabular}
\captionof{table}{Criterios de evaluación. Fuente: elaboración propia.}
\end{center}

\vspace{0.3em}
\noindent\textit{La calidad del informe, la profundidad del análisis y la
correcta interpretación de los resultados impactan de forma transversal en
todas las partes.}

\noindent\rule{\linewidth}{0.4pt}

% =============================================================================
\section*{Notas Importantes}
% =============================================================================

\begin{itemize}[topsep=2pt]
    \item Pueden usar software para derivar y para encontrar los puntos
          críticos reales, o hacerlo a mano. En cualquier caso, el
          procedimiento debe quedar plasmado en el informe.
    \item Usen los \textbf{mismos puntos iniciales} en las Partes~B, C y D para
          que la comparación de la Parte~E sea válida.
    \item El gradiente descendente es sensible al paso $\alpha$: si el método
          diverge, reduzca $\alpha$; si es muy lento, auméntelo con cuidado.
    \item Newton-Raphson requiere que el Hessiano sea invertible en cada
          iteración. Si encuentra problemas numéricos, repórtelos en el
          informe.
    \item En el informe y en el código deben figurar los \textbf{nombres de
          todos los integrantes} del grupo.
\end{itemize}

\vspace{1em}
\begin{center}
\textit{Consultas al correo
\href{mailto:vramireza@udd.cl}{\texttt{vramireza@udd.cl}} o al
\texttt{+56\,9\,9912\,1814}
(respondo siempre que puedo, dentro de horarios razonables).}
\end{center}

% =============================================================================
\end{document}
% =============================================================================
